\section{THIẾT KẾ CƠ SỞ DỮ LIỆU}
\label{sec:chuong3}

\subsection{Các quy tắc ánh xạ đã áp dụng}
\label{subsec:anhxa}

Nhóm chuyển mô hình thực thể liên kết sang mô hình dữ liệu quan hệ theo hai quy tắc chính \cite{elmasri,silberschatz}. Thứ nhất, mỗi thực thể mạnh trở thành một bảng (quan hệ), các thuộc tính đơn trở thành các cột và thuộc tính khóa trở thành khóa chính (primary key); những thuộc tính cũng có giá trị duy nhất như ten\_to hay ten\_dang\_nhap được giữ làm khóa dự tuyển bằng ràng buộc UNIQUE. Thứ hai, với mỗi liên kết 1 -- N, nhóm đưa khóa chính của bảng phía 1 sang bảng phía N làm khóa ngoại (foreign key), nên cả 7 liên kết đều không cần tạo thêm bảng trung gian. Kết quả là 7 bảng, 35 cột, 7 khóa chính và 7 khóa ngoại.

Riêng hai liên kết Đợt cũ và Đợt mới cùng nối Quá trình cư trú với Biến động cư trú, nên bảng bien\_dong\_cu\_tru nhận hai khóa ngoại cùng trỏ về qua\_trinh\_cu\_tru và nhóm đặt tên khác nhau là ma\_qua\_trinh\_cu và ma\_qua\_trinh\_moi để phân biệt vai trò. Hai cột này cho phép NULL vì một biến động có thể chỉ có đợt cũ hoặc chỉ có đợt mới, trong khi các khóa ngoại còn lại đều là NOT NULL do phía 1 của liên kết là bắt buộc.

Mọi khóa chính đều dùng kiểu BIGINT GENERATED BY DEFAULT AS IDENTITY. Identity hỗ trợ cấp mã tự động; ràng buộc PRIMARY KEY bảo đảm tính duy nhất. Vì BY DEFAULT vẫn cho phép nhập mã tường minh, khi nạp dữ liệu có sẵn cần đồng bộ sequence để lần cấp mã tiếp theo không xung đột \cite{pgidentity}. Các giá trị liệt kê như loại cư trú hay loại biến động được lưu bằng mã không dấu (THUONG\_TRU, CHUYEN\_NOI\_BO\ldots) để so sánh ổn định, không phụ thuộc cách gõ dấu tiếng Việt.

\subsection{Lược đồ quan hệ}
\label{subsec:luocdo}

Trong lược đồ, gạch chân biểu thị khóa chính và chữ nghiêng biểu thị khóa ngoại.

{\sloppy
\begin{itemize}[leftmargin=*]
    \item to\_dan\_pho(\underline{ma\_to}, ten\_to, ghi\_chu)
    \item dia\_chi(\underline{ma\_dia\_chi}, \textit{ma\_to}, dia\_chi\_chuan, ghi\_chu)
    \item ho\_cu\_tru(\underline{ma\_ho}, \textit{ma\_dia\_chi}, ngay\_lap, ghi\_chu)
    \item cu\_dan(\underline{ma\_cu\_dan}, ho\_ten, ngay\_sinh, gioi\_tinh, ghi\_chu)
    \item qua\_trinh\_cu\_tru(\underline{ma\_qua\_trinh}, \textit{ma\_cu\_dan}, \textit{ma\_ho}, loai\_cu\_tru, ngay\_bat\_dau, ngay\_ket\_thuc)
    \item tai\_khoan(\underline{ma\_tai\_khoan}, ten\_dang\_nhap, ten\_hien\_thi, vai\_tro, trang\_thai)
    \item bien\_dong\_cu\_tru(\underline{ma\_bien\_dong}, loai\_bien\_dong, \textit{ma\_qua\_trinh\_cu}, \textit{ma\_qua\_trinh\_moi}, ngay\_bien\_dong, \textit{ma\_tai\_khoan}, thoi\_diem\_ghi, ghi\_chu)
\end{itemize}
\par}

\subsection{Cấu trúc các bảng}
\label{subsec:cautruc}

Nhóm cài đặt cơ sở dữ liệu trên PostgreSQL phiên bản 17. Các cột tên và địa chỉ dùng kiểu \texttt{VARCHAR}. Các cột ghi chú dùng kiểu \texttt{TEXT}. Mã hóa UTF-8 hỗ trợ lưu tiếng Việt. Các bảng \ref{tab:todanpho} đến \ref{tab:biendong} mô tả chi tiết từng bảng.

% Định dạng chung cho bảng cấu trúc
\newcommand{\bangcot}{\begin{tabularx}{\textwidth}{|P{3.3cm}|P{3.3cm}|P{3.5cm}|Y|}}
\newcommand{\dongtieude}{\hline \textbf{Cột} & \textbf{Kiểu dữ liệu} & \textbf{Ràng buộc} & \textbf{Mô tả} \\ \hline}

\begin{table}[H]
\centering
\fontsize{12pt}{16pt}\selectfont
\renewcommand{\arraystretch}{1.15}
\bangcot
\dongtieude
ma\_to & BIGINT & PK, tự tăng & Mã tổ dân phố \\ \hline
ten\_to & VARCHAR(60) & NOT NULL, UNIQUE, không rỗng & Tên tổ, không trùng nhau \\ \hline
ghi\_chu & TEXT & NULL & Ghi chú \\ \hline
\end{tabularx}
\caption{Cấu trúc bảng to\_dan\_pho -- lưu danh sách tổ dân phố}
\label{tab:todanpho}
\end{table}

\begin{table}[H]
\centering
\fontsize{12pt}{16pt}\selectfont
\renewcommand{\arraystretch}{1.15}
\bangcot
\dongtieude
ma\_dia\_chi & BIGINT & PK, tự tăng & Mã địa chỉ \\ \hline
ma\_to & BIGINT & NOT NULL, FK $\rightarrow$ to\_dan\_pho & Tổ quản lý địa chỉ \\ \hline
dia\_chi\_chuan & VARCHAR(200) & NOT NULL, không rỗng; UNIQUE cùng ma\_to & Địa chỉ ghi theo dạng chuẩn \\ \hline
ghi\_chu & TEXT & NULL & Ghi chú \\ \hline
\end{tabularx}
\caption{Cấu trúc bảng dia\_chi -- lưu các địa chỉ thuộc từng tổ}
\label{tab:diachi}
\end{table}

\begin{table}[H]
\centering
\fontsize{12pt}{16pt}\selectfont
\renewcommand{\arraystretch}{1.15}
\bangcot
\dongtieude
ma\_ho & BIGINT & PK, tự tăng & Mã hộ \\ \hline
ma\_dia\_chi & BIGINT & NOT NULL, FK $\rightarrow$ dia\_chi & Địa chỉ nơi hộ sinh sống \\ \hline
ngay\_lap & DATE & NOT NULL & Ngày lập hộ \\ \hline
ghi\_chu & TEXT & NULL & Ghi chú \\ \hline
\end{tabularx}
\caption{Cấu trúc bảng ho\_cu\_tru -- lưu các hộ cư trú}
\label{tab:hocutru}
\end{table}

\begin{table}[H]
\centering
\fontsize{12pt}{16pt}\selectfont
\renewcommand{\arraystretch}{1.15}
\bangcot
\dongtieude
ma\_cu\_dan & BIGINT & PK, tự tăng & Mã cư dân \\ \hline
ho\_ten & VARCHAR(120) & NOT NULL, không rỗng & Họ và tên \\ \hline
ngay\_sinh & DATE & NOT NULL & Ngày sinh \\ \hline
gioi\_tinh & VARCHAR(10) & NOT NULL, CHECK (Nam, Nu, Khac) & Giới tính \\ \hline
ghi\_chu & TEXT & NULL & Ghi chú \\ \hline
\end{tabularx}
\caption{Cấu trúc bảng cu\_dan -- lưu thông tin cá nhân của cư dân}
\label{tab:cudan}
\end{table}

\begin{table}[H]
\centering
\fontsize{12pt}{16pt}\selectfont
\renewcommand{\arraystretch}{1.15}
\bangcot
\dongtieude
ma\_qua\_trinh & BIGINT & PK, tự tăng & Mã quá trình cư trú \\ \hline
ma\_cu\_dan & BIGINT & NOT NULL, FK $\rightarrow$ cu\_dan & Cư dân của đợt cư trú \\ \hline
ma\_ho & BIGINT & NOT NULL, FK $\rightarrow$ ho\_cu\_tru & Hộ mà cư dân cư trú \\ \hline
loai\_cu\_tru & VARCHAR(20) & NOT NULL, CHECK (THUONG\_TRU, TAM\_TRU) & Loại cư trú \\ \hline
ngay\_bat\_dau & DATE & NOT NULL & Ngày bắt đầu đợt cư trú \\ \hline
ngay\_ket\_thuc & DATE & NULL, $\geq$ ngay\_bat\_dau & Ngày kết thúc; NULL nghĩa là chưa kết thúc \\ \hline
\end{tabularx}
\caption{Cấu trúc bảng qua\_trinh\_cu\_tru -- lưu từng đợt cư trú của cư dân trong hộ}
\label{tab:quatrinh}
\end{table}

\begin{table}[H]
\centering
\fontsize{12pt}{16pt}\selectfont
\renewcommand{\arraystretch}{1.15}
\bangcot
\dongtieude
ma\_tai\_khoan & BIGINT & PK, tự tăng & Mã tài khoản \\ \hline
ten\_dang\_nhap & VARCHAR(60) & NOT NULL, UNIQUE, không rỗng & Tên đăng nhập \\ \hline
ten\_hien\_thi & VARCHAR(120) & NOT NULL, không rỗng & Tên hiển thị của cán bộ \\ \hline
vai\_tro & VARCHAR(20) & NOT NULL, CHECK (QUAN\_TRI, CAN\_BO, THONG\_KE) & Vai trò nghiệp vụ của người ghi \\ \hline
trang\_thai & VARCHAR(20) & NOT NULL, mặc định HOAT\_DONG, CHECK (HOAT\_DONG, KHOA) & Trạng thái tài khoản \\ \hline
\end{tabularx}
\caption{Cấu trúc bảng tai\_khoan -- lưu tài khoản cán bộ sử dụng hệ thống}
\label{tab:taikhoan}
\end{table}

\begin{table}[H]
\centering
\fontsize{12pt}{16pt}\selectfont
\renewcommand{\arraystretch}{1.15}
\bangcot
\dongtieude
ma\_bien\_dong & BIGINT & PK, tự tăng & Mã biến động \\ \hline
loai\_bien\_dong & VARCHAR(20) & NOT NULL & CHUYEN\_DEN, CHUYEN\_DI hoặc CHUYEN\_NOI\_BO \\ \hline
ma\_qua\_trinh\_cu & BIGINT & NULL, FK $\rightarrow$ qua\_trinh\_cu\_tru & Đợt cư trú bị kết thúc \\ \hline
ma\_qua\_trinh\_moi & BIGINT & NULL, FK $\rightarrow$ qua\_trinh\_cu\_tru & Đợt cư trú được mở \\ \hline
ngay\_bien\_dong & DATE & NOT NULL & Ngày biến động xảy ra \\ \hline
ma\_tai\_khoan & BIGINT & NOT NULL, FK $\rightarrow$ tai\_khoan & Tài khoản ghi nhận \\ \hline
thoi\_diem\_ghi & TIMESTAMPTZ & NOT NULL, mặc định CURRENT\_TIMESTAMP & Thời điểm nhập vào hệ thống \\ \hline
ghi\_chu & TEXT & NULL & Ghi chú \\ \hline
\end{tabularx}
\caption{Cấu trúc bảng bien\_dong\_cu\_tru -- lưu lịch sử biến động cư trú}
\label{tab:biendong}
\end{table}

\subsection{Khóa, phụ thuộc hàm và chuẩn hóa}
\label{subsec:chuanhoa}

Nhóm phân tích chuẩn hóa dựa trên khóa ứng viên và phụ thuộc hàm (FD) của từng bảng \cite{elmasri,silberschatz}. Mỗi khóa chính xác định mọi thuộc tính còn lại trong bảng. Ngoài ra, \texttt{ten\_to} là khóa ứng viên của bảng tổ; cặp (\texttt{ma\_to}, \texttt{dia\_chi\_chuan}) là khóa ứng viên của bảng địa chỉ; \texttt{ten\_dang\_nhap} là khóa ứng viên của bảng tài khoản. Các khóa này đều được bảo đảm bởi NOT NULL và UNIQUE. B06 còn bảo đảm bộ (\texttt{ma\_cu\_dan}, \texttt{loai\_cu\_tru}, \texttt{ngay\_bat\_dau}) định danh một đợt: hai đợt cùng người, cùng loại và cùng ngày bắt đầu sẽ chồng lấn nên bị từ chối.

Ví dụ, \texttt{ma\_dia\_chi} xác định \texttt{ma\_to}, \texttt{dia\_chi\_chuan} và \texttt{ghi\_chu}; cặp tổ và địa chỉ chuẩn cũng xác định mã địa chỉ cùng ghi chú. \texttt{ma\_ho} xác định địa chỉ và ngày lập hộ; \texttt{ma\_cu\_dan} xác định thông tin hồ sơ; \texttt{ma\_qua\_trinh} xác định cư dân, hộ, loại và khoảng ngày của đợt. Tên tổ, địa chỉ dạng chữ và tên tài khoản được lưu ở bảng tương ứng, không sao chép vào bảng đợt hoặc sự kiện.

\begin{itemize}[leftmargin=*]
    \item \textbf{1NF:} mỗi ô chứa một giá trị, không lưu danh sách thành viên trong một cột; các lần cư trú được tách thành các dòng riêng. Với thuộc tính tùy chọn, CSDL dùng NULL theo ngữ nghĩa SQL; khi vận dụng đại số quan hệ cổ điển cần chuẩn bị đầu vào như mục \ref{subsec:qd}.
    \item \textbf{2NF:} xét cả khóa chính lẫn khóa ứng viên ghép. Với địa chỉ, mã và ghi chú của bản ghi phụ thuộc đầy đủ vào cặp tổ và địa chỉ chuẩn. Với đợt cư trú, không có FD được công bố từ một phần thực sự của khóa ghép nêu trên tới thuộc tính không khóa; chẳng hạn một cư dân hoặc một loại cư trú riêng lẻ không xác định hộ hay ngày kết thúc.
    \item \textbf{3NF:} theo tập FD đã công bố, sáu bảng tổ, địa chỉ, hộ, cư dân, quá trình cư trú và tài khoản không có FD không tầm thường với vế trái không là siêu khóa và vế phải là thuộc tính không khóa. Việc tách tên tổ, địa chỉ và thông tin tài khoản tránh cập nhật lặp lại các giá trị này ở nhiều bản ghi lịch sử.
\end{itemize}

\textbf{Ngày biến động và chuẩn hóa.} Xét các phụ thuộc hàm do khóa chính xác định, bảng biến động đáp ứng 3NF. B09 tạo thêm phụ thuộc có điều kiện trên các dòng có khóa ngoại đợt khác NULL. Nếu biến động có đợt mới, \texttt{ngay\_bien\_dong} bằng ngày bắt đầu đợt mới. Nếu có đợt cũ, ngày biến động bằng ngày kết thúc đợt cũ cộng một. Các khóa ngoại này có thể lặp, nên mã đợt xác định ngày biến động trong từng vai trò mà vẫn có thể tham chiếu từ nhiều bản ghi.

Nhóm lưu ngày biến động ngay trong bản ghi để tra cứu ngày nghiệp vụ. Trigger B09 kiểm tra sự nhất quán theo cả hai chiều. Phép kiểm E1 phát hiện một đợt được dùng lặp trong cùng vai trò.

\subsection{Các ràng buộc dữ liệu}
\label{subsec:rangbuoc}

Nhóm chia ràng buộc thành hai lớp. Lớp thứ nhất nằm ngay trong câu lệnh tạo bảng (tệp \texttt{sql/01\_schema.sql}), gồm khóa chính, khóa ngoại, UNIQUE và CHECK trên một dòng. Lớp thứ hai là các quy tắc liên quan tới nhiều dòng hoặc nhiều bảng mà CHECK không diễn đạt được, nhóm cài trong tệp \texttt{sql/02\_constraints.sql} bằng ràng buộc loại trừ (exclusion constraint), trigger và hàm PL/pgSQL.

\begin{enumerate}[leftmargin=*]
    \item \textbf{Miền giá trị (CHECK):} giới tính chỉ nhận Nam, Nu hoặc Khac; loại cư trú chỉ nhận THUONG\_TRU hoặc TAM\_TRU; vai trò và trạng thái tài khoản chỉ nhận các mã đã liệt kê; ngày kết thúc rỗng hoặc không sớm hơn ngày bắt đầu (B04). Tên tổ, địa chỉ, họ tên, tên đăng nhập và tên hiển thị không được rỗng hoặc chỉ toàn khoảng trắng (B15); năm ràng buộc CHECK này (\texttt{ck\_b15\_*}) được bổ sung trong \texttt{02\_constraints.sql}, nên sau khi chạy đủ hai tệp cơ sở dữ liệu có tổng cộng 11 ràng buộc CHECK.
    \item \textbf{Hình dạng biến động (B07):} ràng buộc \texttt{ck\_bd\_hinh\_dang} yêu cầu CHUYEN\_DEN chỉ có đợt mới, CHUYEN\_DI chỉ có đợt cũ, còn CHUYEN\_NOI\_BO có cả hai đợt và hai đợt này khác nhau.
    \item \textbf{Không chồng lấn đợt cư trú (B06):} nhóm dùng ràng buộc loại trừ \texttt{ex\_qt\_cu\_tru\_khong\_chong\_lan} với chỉ mục GiST (cần extension \texttt{btree\_gist}). Hai đợt cùng ma\_cu\_dan, cùng loai\_cu\_tru sẽ bị từ chối nếu khoảng ngày \texttt{daterange(ngay\_bat\_dau, ngay\_ket\_thuc, '[]')} của chúng giao nhau (toán tử \texttt{\&\&}); ngày kết thúc NULL được hiểu là không giới hạn.
    \item \textbf{Khớp sự kiện với đợt (B08, B09):} trigger \texttt{trg\_bien\_dong\_hop\_le} kiểm tra hai đợt của chuyển nội bộ cùng cư dân, khác hộ, đợt mới bắt đầu đúng ngày chuyển và đợt cũ kết thúc đúng ngày trước đó. Trigger \texttt{trg\_dot\_khong\_pha\_su\_kien} kiểm tra chiều ngược lại: sửa ngày của một đợt không được làm sai biến động đã ghi.
    \item \textbf{Giao dịch nguyên tử (B10):} hàm \texttt{fn\_chuyen\_noi\_bo} đóng đợt cũ, mở đợt mới và ghi biến động trong cùng một lần gọi; nếu một bước lỗi thì toàn bộ bị hủy.
    \item \textbf{Giữ địa chỉ lịch sử (B11):} trigger không cho đổi địa chỉ của một hộ, hay đổi tổ của một địa chỉ, khi đã có lịch sử cư trú tham chiếu tới.
    \item \textbf{Không xóa lịch sử (B12):} mọi khóa ngoại đều không CASCADE nên không xóa được hồ sơ đang được tham chiếu; ngoài ra role nghiệp vụ csdl17\_admin bị thu hồi quyền DELETE trên hai bảng lịch sử. Phạm vi quyền cập nhật được trình bày tại mục \ref{subsec:danhgia}.
    \item \textbf{Tài khoản hoạt động (B14):} trigger chỉ cho ghi biến động mới bằng tài khoản có trạng thái HOAT\_DONG; khóa tài khoản về sau không làm mất các biến động cũ.
    \item \textbf{Ngày hợp lý (B15):} trigger kiểm tra ngày sinh và ngày lập hộ không muộn hơn ngày bắt đầu các đợt cư trú, theo cả chiều ghi đợt và chiều sửa hồ sơ.
\end{enumerate}

Phép kiểm E1 trong \texttt{06\_validation.sql} phát hiện các biến động dùng lặp một đợt trong cùng vai trò. Điểm này cùng các giới hạn X01 và X06 được tổng hợp trong phần Kết luận.

Nhóm đã dựng lại cơ sở dữ liệu từ đầu trên PostgreSQL 17.12, chạy lần lượt \texttt{01\_schema.sql}, \texttt{02\_constraints.sql}, \texttt{03\_seed.sql} và \texttt{06\_validation.sql}; bước kiểm tra dữ liệu báo không có vi phạm B06, B08, B09, E1 và B15. Bộ dữ liệu mẫu nạp được 636 bản ghi: 8 tổ, 50 địa chỉ, 60 hộ, 200 cư dân, 250 đợt cư trú, 5 tài khoản và 63 biến động. Bảng \ref{tab:thurangbuoc} là kết quả thử một số câu lệnh cố tình vi phạm ràng buộc; mỗi lệnh chạy trong giao dịch được hoàn tác, và cột bên phải ghi đúng mã lỗi (SQLSTATE) cùng tên ràng buộc hoặc thông báo mà PostgreSQL trả về.

\begin{table}[H]
\centering
\fontsize{12pt}{16pt}\selectfont
\renewcommand{\arraystretch}{1.15}
\begin{tabularx}{\textwidth}{|Y|P{6.4cm}|}
\hline
\textbf{Thao tác thử} & \textbf{Kết quả của PostgreSQL} \\
\hline
Mở thêm đợt THUONG\_TRU cho cư dân 1 trùng thời gian với đợt đang có & Từ chối (23P01): vi phạm \texttt{ex\_qt\_cu\_tru\_khong\_chong\_lan} (B06) \\
\hline
Ghi CHUYEN\_DEN nhưng có cả đợt cũ: (a) đợt cũ là đợt 1 đang mở; (b) đợt cũ 16 kết thúc đúng ngày trước ngày chuyển, đợt mới 17 bắt đầu đúng ngày chuyển & Từ chối (23514). (a) Trigger B09 chặn trước: ``đợt cũ 1 phải kết thúc trước ngày chuyển 2026-10-15 đúng 1 ngày''. (b) Ngày khớp nên trigger cho qua, CHECK \texttt{ck\_bd\_hinh\_dang} (B07) chặn \\
\hline
Đổi địa chỉ của hộ 1 đã có lịch sử & Từ chối (23514): ``B11: không đổi địa chỉ của hộ 1 đã có lịch sử'' \\
\hline
Đợt có ngày kết thúc trước ngày bắt đầu & Từ chối (23514): vi phạm \texttt{ck\_qt\_ngay} \\
\hline
Xóa cư dân 1 đang có đợt cư trú & Từ chối (23503): vi phạm khóa ngoại \texttt{qua\_trinh\_cu\_tru\_ma\_cu\_dan\_fkey} \\
\hline
Thêm cư dân có họ tên chỉ gồm khoảng trắng & Từ chối (23514): vi phạm \texttt{ck\_b15\_ho\_ten} \\
\hline
\end{tabularx}
\caption{Kết quả thử các câu lệnh vi phạm ràng buộc}
\label{tab:thurangbuoc}
\end{table}

\subsection{Script tạo cơ sở dữ liệu}
\label{subsec:script}

Đoạn mã \ref{lst:schema} là phần tạo bảng trong \texttt{sql/01\_schema.sql}, kèm ràng buộc B06 lấy từ \texttt{sql/02\_constraints.sql}. Phụ lục~\hyperref[sec:phuluc]{A} trình bày mã nguồn các trigger và hàm tiêu biểu. Phần tạo bảng được chép nguyên văn từ tệp PostgreSQL của dự án, chỉ lược bỏ mệnh đề \texttt{IF NOT EXISTS} và năm câu lệnh tạo chỉ mục; ràng buộc B06 trong tệp gốc được bọc trong một khối \texttt{DO} kiểm tra tên ràng buộc để có thể chạy lại script nhiều lần.

\begin{lstlisting}[style=sqlvn, caption={Script tạo bảng và ràng buộc B06 trên PostgreSQL}, label={lst:schema}]
CREATE TABLE to_dan_pho (
    ma_to          BIGINT GENERATED BY DEFAULT AS IDENTITY PRIMARY KEY,
    ten_to         VARCHAR(60) NOT NULL UNIQUE,
    ghi_chu        TEXT
);

CREATE TABLE dia_chi (
    ma_dia_chi     BIGINT GENERATED BY DEFAULT AS IDENTITY PRIMARY KEY,
    ma_to          BIGINT NOT NULL REFERENCES to_dan_pho(ma_to),
    dia_chi_chuan  VARCHAR(200) NOT NULL,
    ghi_chu        TEXT,
    CONSTRAINT uq_dia_chi_to UNIQUE (ma_to, dia_chi_chuan)
);

CREATE TABLE ho_cu_tru (
    ma_ho          BIGINT GENERATED BY DEFAULT AS IDENTITY PRIMARY KEY,
    ma_dia_chi     BIGINT NOT NULL REFERENCES dia_chi(ma_dia_chi),
    ngay_lap       DATE NOT NULL,
    ghi_chu        TEXT
);

CREATE TABLE cu_dan (
    ma_cu_dan      BIGINT GENERATED BY DEFAULT AS IDENTITY PRIMARY KEY,
    ho_ten         VARCHAR(120) NOT NULL,
    ngay_sinh      DATE NOT NULL,
    gioi_tinh      VARCHAR(10) NOT NULL,
    ghi_chu        TEXT,
    CONSTRAINT ck_cu_dan_gioi_tinh CHECK (gioi_tinh IN ('Nam', 'Nu', 'Khac'))
);

CREATE TABLE qua_trinh_cu_tru (
    ma_qua_trinh   BIGINT GENERATED BY DEFAULT AS IDENTITY PRIMARY KEY,
    ma_cu_dan      BIGINT NOT NULL REFERENCES cu_dan(ma_cu_dan),
    ma_ho          BIGINT NOT NULL REFERENCES ho_cu_tru(ma_ho),
    loai_cu_tru    VARCHAR(20) NOT NULL,
    ngay_bat_dau   DATE NOT NULL,
    ngay_ket_thuc  DATE,
    CONSTRAINT ck_qt_loai CHECK (loai_cu_tru IN ('THUONG_TRU', 'TAM_TRU')),
    CONSTRAINT ck_qt_ngay CHECK (ngay_ket_thuc IS NULL OR ngay_ket_thuc >= ngay_bat_dau)
);

CREATE TABLE tai_khoan (
    ma_tai_khoan   BIGINT GENERATED BY DEFAULT AS IDENTITY PRIMARY KEY,
    ten_dang_nhap  VARCHAR(60) NOT NULL UNIQUE,
    ten_hien_thi   VARCHAR(120) NOT NULL,
    vai_tro        VARCHAR(20) NOT NULL,
    trang_thai     VARCHAR(20) NOT NULL DEFAULT 'HOAT_DONG',
    CONSTRAINT ck_tk_vai_tro CHECK (vai_tro IN ('QUAN_TRI', 'CAN_BO', 'THONG_KE')),
    CONSTRAINT ck_tk_trang_thai CHECK (trang_thai IN ('HOAT_DONG', 'KHOA'))
);

CREATE TABLE bien_dong_cu_tru (
    ma_bien_dong       BIGINT GENERATED BY DEFAULT AS IDENTITY PRIMARY KEY,
    loai_bien_dong     VARCHAR(20) NOT NULL,
    ma_qua_trinh_cu    BIGINT REFERENCES qua_trinh_cu_tru(ma_qua_trinh),
    ma_qua_trinh_moi   BIGINT REFERENCES qua_trinh_cu_tru(ma_qua_trinh),
    ngay_bien_dong     DATE NOT NULL,
    ma_tai_khoan       BIGINT NOT NULL REFERENCES tai_khoan(ma_tai_khoan),
    thoi_diem_ghi      TIMESTAMPTZ NOT NULL DEFAULT CURRENT_TIMESTAMP,
    ghi_chu            TEXT,
    CONSTRAINT ck_bd_hinh_dang CHECK (
        (loai_bien_dong = 'CHUYEN_DEN'
         AND ma_qua_trinh_cu IS NULL AND ma_qua_trinh_moi IS NOT NULL)
     OR (loai_bien_dong = 'CHUYEN_DI'
         AND ma_qua_trinh_cu IS NOT NULL AND ma_qua_trinh_moi IS NULL)
     OR (loai_bien_dong = 'CHUYEN_NOI_BO'
         AND ma_qua_trinh_cu IS NOT NULL AND ma_qua_trinh_moi IS NOT NULL
         AND ma_qua_trinh_cu <> ma_qua_trinh_moi)
    )
);

-- B06: cùng người, cùng loại cư trú thì khoảng ngày không được giao nhau
CREATE EXTENSION IF NOT EXISTS btree_gist;
ALTER TABLE qua_trinh_cu_tru
  ADD CONSTRAINT ex_qt_cu_tru_khong_chong_lan
  EXCLUDE USING gist (
    ma_cu_dan WITH =,
    loai_cu_tru WITH =,
    daterange(ngay_bat_dau, ngay_ket_thuc, '[]') WITH &&
  );
\end{lstlisting}

\subsection{Chỉ mục}
\label{subsec:chimuc}

PostgreSQL tự tạo chỉ mục B-tree cho khóa chính và ràng buộc UNIQUE. Nhóm bổ sung năm chỉ mục sau cho các cột thường dùng khi kết nối, lọc hoặc kiểm tra ràng buộc \cite{postgres}.

\begin{lstlisting}[style=sqlvn, caption={Năm chỉ mục bổ sung}, label={lst:chimuc}]
CREATE INDEX IF NOT EXISTS idx_dia_chi_to ON dia_chi(ma_to);
CREATE INDEX IF NOT EXISTS idx_ho_dia_chi ON ho_cu_tru(ma_dia_chi);
CREATE INDEX IF NOT EXISTS idx_qt_cu_dan_ngay ON qua_trinh_cu_tru(ma_cu_dan, ngay_bat_dau);
CREATE INDEX IF NOT EXISTS idx_qt_ho_ngay ON qua_trinh_cu_tru(ma_ho, ngay_bat_dau);
CREATE INDEX IF NOT EXISTS idx_bd_loai_ngay ON bien_dong_cu_tru(loai_bien_dong, ngay_bien_dong);
\end{lstlisting}

\begin{itemize}[leftmargin=*]
    \item \texttt{idx\_dia\_chi\_to} và \texttt{idx\_ho\_dia\_chi} hỗ trợ đường nối từ tổ đến địa chỉ và hộ, chẳng hạn khi thống kê theo tổ (QSQL03) hoặc tìm địa chỉ chưa có hộ (QSQL04).
    \item \texttt{idx\_qt\_cu\_dan\_ngay} hỗ trợ tra cứu đợt theo cư dân và khoảng ngày, chẳng hạn khi kiểm tra đợt còn hiệu lực ở QSQL07 hoặc khi trigger B15 kiểm tra ngày sinh.
    \item \texttt{idx\_qt\_ho\_ngay} hỗ trợ tìm các đợt của một hộ, chẳng hạn khi đếm người ở QSQL02 hoặc khi kiểm tra lịch sử của hộ theo B11.
    \item \texttt{idx\_bd\_loai\_ngay} phục vụ truy vấn lọc theo loại biến động và, khi cần, theo khoảng ngày. QSQL09 lọc chuyển nội bộ và sắp theo mã biến động. QSQL08 tổng hợp toàn bộ các loại theo tháng, nên chỉ mục này không bảo đảm truy vấn tránh quét bảng.
\end{itemize}

Với chỉ mục B-tree nhiều cột, đặt cột so sánh bằng ở đầu hỗ trợ tìm nhanh phần cần đọc, sau đó có thể giới hạn trên cột ngày. Lọc theo khoảng ngày tạo điều kiện dùng chỉ mục trên cột ngày. PostgreSQL cũng hỗ trợ chỉ mục trên biểu thức \cite{pgexpressionindex}. Trên bộ 636 bản ghi, lựa chọn kế hoạch phụ thuộc số dòng cần đọc và chi phí ước tính. Nhóm chưa đo mức cải thiện thời gian chạy.

\subsection{Câu lệnh DML minh họa}
\label{subsec:dml}

Ba ví dụ minh họa cách các ràng buộc ở mục \ref{subsec:rangbuoc} kiểm tra thao tác thêm và sửa dữ liệu. Nhóm chạy trên PostgreSQL 17.12 với cơ sở dữ liệu dựng mới từ schema, ràng buộc và bộ 636 bản ghi. Hai ví dụ đầu lấy ngày hệ thống làm ngày chuyển. Mỗi giao dịch được hoàn tác để giữ nguyên các bản ghi mẫu. Script và kết quả được lưu lần lượt tại \texttt{tests/report\_dml\_minh\_hoa.sql} và \texttt{tests/report\_dml\_minh\_hoa\_log.txt}.

\textbf{Ví dụ 1: chuyển nội bộ bằng hàm.} Cư dân 40 (Phạm Thu Hạnh) đang thường trú tại hộ 21 theo đợt 44. Cán bộ dùng tài khoản 2 ghi nhận người này chuyển thường trú sang hộ 2 từ hôm nay.

\begin{lstlisting}[style=sqlvn]
BEGIN;
SELECT * FROM fn_chuyen_noi_bo(44, 2, 'THUONG_TRU', CURRENT_DATE, 2);
ROLLBACK;
\end{lstlisting}

Hàm khóa đợt 44 bằng \texttt{FOR UPDATE}, đặt ngày kết thúc của đợt này là hôm qua, thêm đợt thường trú mới tại hộ 2 bắt đầu từ hôm nay, rồi ghi một biến động CHUYEN\_NOI\_BO trỏ tới cả hai đợt. Kết quả trả về là mã đợt mới và mã biến động vừa được cấp; trên cơ sở dữ liệu dựng mới, đó là đợt 251 và biến động 64. Khi biến động được ghi, trigger \texttt{trg\_bien\_dong\_hop\_le} vẫn kiểm tra lại B08, B09 và B14; nếu một bước lỗi thì cả ba thay đổi bị hủy (B10). Người dùng chỉ cần một lệnh, không phải tự viết ba câu UPDATE và INSERT theo đúng thứ tự.

\textbf{Ví dụ 2: ghi chuyển đi.} Cư dân 44 (Đặng Ngọc Chi) đang thường trú tại hộ 16 theo đợt 50 và chuyển ra khỏi phường từ hôm nay. Nhóm không viết hàm riêng cho trường hợp này, nên hai lệnh được đặt trong cùng một giao dịch.

\begin{lstlisting}[style=sqlvn]
BEGIN;
UPDATE qua_trinh_cu_tru
SET ngay_ket_thuc = CURRENT_DATE - 1
WHERE ma_qua_trinh = 50;

INSERT INTO bien_dong_cu_tru (loai_bien_dong, ma_qua_trinh_cu, ma_qua_trinh_moi,
                              ngay_bien_dong, ma_tai_khoan)
VALUES ('CHUYEN_DI', 50, NULL, CURRENT_DATE, 2);
ROLLBACK;
\end{lstlisting}

Thứ tự hai lệnh là bắt buộc. Nếu INSERT chạy trước khi đợt 50 được đóng, trigger B09 từ chối với lỗi 23514 ``B09: đợt cũ 50 phải kết thúc trước ngày chuyển \ldots{} đúng 1 ngày'', trong đó ngày chuyển là ngày chạy lệnh. Biến động CHUYEN\_DI chỉ có đợt cũ, đúng hình dạng mà CHECK \texttt{ck\_bd\_hinh\_dang} yêu cầu (B07).

\textbf{Ví dụ 3: UPDATE bị trigger chặn.} Đợt 17 là đợt mới của biến động 2, lần chuyển nội bộ ngày 15/10/2026 của cư dân 15. Câu lệnh sau thử dời ngày bắt đầu của đợt này sang 20/10/2026.

\begin{lstlisting}[style=sqlvn]
UPDATE qua_trinh_cu_tru
SET ngay_bat_dau = DATE '2026-10-20'
WHERE ma_qua_trinh = 17;
\end{lstlisting}

Trigger \texttt{trg\_dot\_khong\_pha\_su\_kien} kiểm tra lại các biến động đang trỏ tới đợt 17 và báo lỗi 23514: ``B09: đợt mới 17 phải bắt đầu đúng ngày chuyển 2026-10-15''. Dòng dữ liệu giữ nguyên, nên biến động 2 vẫn khớp với hai đợt mà nó trỏ tới.
